% !TeX program = lualatex
% halo:begin
% {
%   "version": 1,
%   "publish": false,
%   "course": "mathematical-analysis-i",
%   "section": "6.4",
%   "title": "数学分析（一）6.4：函数的极值与最大（小）值",
%   "categories": ["study-notes"],
%   "tags": ["mathematics"],
%   "excerpt": "用一阶导数的变号与首个非零高阶导数判别极值，再用端点比较或全区间单调性确定最值。"
% }
% halo:end
% 来源：IMG_3596、IMG_3597 左页；华东师范大学《数学分析》第5版上册，第六章§4，书页132–137（PDF 143–148），含例2–6、习题6.4第2、5题。
\RequirePackage{iftex}
\RequireLuaTeX
\documentclass[UTF8,fontset=fandol,a4paper,12pt]{ctexart}
\usepackage[margin=2.5cm]{geometry}
\usepackage{amsmath,amssymb,amsthm}
\usepackage{fancyhdr}
\usepackage[hidelinks]{hyperref}
\theoremstyle{definition}
\newtheorem{theorem}{定理}[section]
\newtheorem{proposition}[theorem]{命题}
\newtheorem{example}[theorem]{例题}
\renewcommand{\proofname}{证明}
\setlength{\headheight}{16pt}
\pagestyle{fancy}\fancyhf{}
\fancyhead[L]{\small 数学分析（一）}
\fancyhead[R]{\small 6.4\quad 极值与最大（小）值}
\fancyfoot[C]{\thepage}\renewcommand{\headrulewidth}{0.3pt}
\raggedbottom
\title{函数的极值与最大（小）值}
\author{Yuchen Zhang}
\date{整理日期：2026年11月26日}
\makeatletter
\renewcommand{\maketitle}{\begin{center}{\LARGE\bfseries\@title\par}\vspace{0.6em}\@author\par\vspace{0.2em}{\small\@date\par}\end{center}}
\makeatother
\begin{document}
\maketitle

\section{从两侧的增减判断极值}
极值是局部概念，比较的是函数在某一点附近的值；最值则是整体概念，比较的是函数在整个指定区间上的值。设 $x_0$ 为定义域的内点，若在 $x_0$ 的某个邻域内恒有 $f(x)\ge f(x_0)$，则称 $x_0$ 为 $f$ 的\textbf{极小值点}，$f(x_0)$ 为\textbf{极小值}；若在相应的去心邻域内恒有严格不等式 $f(x)>f(x_0)$，则称为严格极小值点与严格极小值。将不等号反向，即得极大值点与极大值的定义。

当 $f$ 在内点 $x_0$ 处可导且取得极值时，费马定理表明必有 $f'(x_0)=0$，满足该等式的点称为\textbf{稳定点}（或驻点）。在寻找极值时，既不能遗漏不可导点（例如 $f(x)=|x|$ 在 $0$ 处不可导，却取得极小值），也不能把驻点直接当作极值点（例如 $f(x)=x^3$ 在 $0$ 处导数为零，但函数在全轴上严格递增，并不取极值）。要确认一个候选点是否为极值点，需要借助充分条件。

最直观的思路，是考察函数在候选点左右两侧的增减变化。

\begin{theorem}[极值的第一充分条件，教材定理 6.11]
设 $f$ 在 $x_0$ 处连续，并在 $x_0$ 的某去心邻域内可导。若在 $x_0$ 左侧恒有 $f'\le0$、在右侧恒有 $f'\ge0$，则 $f(x_0)$ 为极小值；若在左侧恒有 $f'\ge0$、在右侧恒有 $f'\le0$，则 $f(x_0)$ 为极大值。当两侧均取严格不等号时，所得极值也是严格的。
\end{theorem}

\begin{proof}
以极小值情形为例。在去心邻域的左侧任取 $x<x_0$，由于 $f$ 在 $[x,x_0]$ 上连续、在 $(x,x_0)$ 内可导，由拉格朗日中值定理存在 $\xi\in(x,x_0)$ 使得
\[
 f(x_0)-f(x)=f'(\xi)(x_0-x)\le0.
\]
同理，在右侧任取 $x>x_0$，存在 $\eta\in(x_0,x)$ 使得
\[
 f(x)-f(x_0)=f'(\eta)(x-x_0)\ge0.
\]
合并两侧即知在该邻域内恒有 $f(x)\ge f(x_0)$，故 $f(x_0)$ 为极小值（当两侧导数严格保号时不等式也严格）。极大值的情形对 $-f$ 应用同一论证即可。
\end{proof}

\begin{example}
求 $f(x)=(2x-5)\sqrt[3]{x^2}$ 的极值。
\end{example}

\begin{proof}[解]
函数在 $\mathbb R$ 上处处连续。当 $x\ne0$ 时，求导得
\[
 f'(x)=\frac{10(x-1)}{3\sqrt[3]{x}},
\]
由此得到唯一驻点 $x=1$。此外还需检查 $x=0$，此处差商为 $(2x-5)/\sqrt[3]{x}$，当 $x\to0$ 时无有限极限，故 $x=0$ 为不可导点。按这两个分界点列表考察 $f'$ 的符号变化：
\[
\begin{array}{c|ccccc}
x &(-\infty,0)&0&(0,1)&1&(1,+\infty)\\ \hline
f'(x)&+&\text{不存在}&-&0&+\\
f(x)&\nearrow&0&\searrow&-3&\nearrow
\end{array}
\]
由第一充分条件可知，$x=0$ 是严格极大值点，极大值为 $f(0)=0$；$x=1$ 是严格极小值点，极小值为 $f(1)=-3$。
\end{proof}

\section{从一点的导数判断极值}
第一充分条件需要分析候选点两侧邻域内一阶导数的符号。如果函数在驻点处具有非零的二阶导数，那么只需借助该点处的局部泰勒展开，便能直接判断邻近函数值的高低。

\begin{theorem}[极值的第二充分条件，教材定理 6.12]
设 $f$ 在 $x_0$ 的某邻域内可导，在 $x_0$ 处二阶可导，且 $f'(x_0)=0$。若 $f''(x_0)>0$，则 $x_0$ 是严格极小值点；若 $f''(x_0)<0$，则 $x_0$ 是严格极大值点。
\end{theorem}

\begin{proof}
令 $h=x-x_0$。由带佩亚诺型余项的二阶泰勒公式及 $f'(x_0)=0$，当 $h\to0$ 时有
\[
 f(x_0+h)-f(x_0)=h^2\left(\frac{f''(x_0)}{2}+o(1)\right).
\]
当非零增量 $h$ 充分小时，$h^2>0$，而括号内的表达式与 $f''(x_0)$ 同号，从而 $f(x_0+h)-f(x_0)$ 与 $f''(x_0)$ 保持同号，结论随之成立。
\end{proof}

\begin{example}
求 $f(x)=x^2+432/x$ 的极值。
\end{example}

\begin{proof}[解]
函数的定义域为 $(-\infty,0)\cup(0,+\infty)$，在两段区间上均可导。计算前两阶导数得
\[
 f'(x)=2x-\frac{432}{x^2}=\frac{2(x^3-216)}{x^2},
 \qquad f''(x)=2+\frac{864}{x^3}.
\]
令 $f'(x)=0$ 得唯一驻点 $x=6$。由于 $f''(6)=6>0$，由第二充分条件可知 $f(6)=108$ 是严格极小值。需要留意，$x=0$ 不在定义域内，没有函数值可供比较，不是极值候选点。
\end{proof}

当驻点处同时满足 $f''(x_0)=0$ 时，第二充分条件失效。此时顺着同一思路，只需继续向高阶展开，找出泰勒公式中首个非零的非常数项即可。

\begin{theorem}[极值的第三充分条件，教材定理 6.13]
设 $f$ 在 $x_0$ 的某邻域内有直到 $n-1$ 阶的导数，在 $x_0$ 处有 $n$ 阶导数（$n\ge2$），且
\[
 f'(x_0)=\cdots=f^{(n-1)}(x_0)=0,\qquad f^{(n)}(x_0)\ne0.
\]
若 $n$ 为偶数，则当 $f^{(n)}(x_0)>0$ 时 $x_0$ 为严格极小值点，当 $f^{(n)}(x_0)<0$ 时 $x_0$ 为严格极大值点；若 $n$ 为奇数，则 $x_0$ 不是极值点。
\end{theorem}

\begin{proof}
令 $h=x-x_0$，由带佩亚诺型余项的 $n$ 阶泰勒公式得
\[
 f(x_0+h)-f(x_0)=h^n\left(\frac{f^{(n)}(x_0)}{n!}+o(1)\right).
\]
在充分小的去心邻域内，括号内的表达式与 $f^{(n)}(x_0)$ 同号。当 $n$ 为偶数时，$h^n$ 在 $x_0$ 两侧均为正，故 $f(x_0+h)-f(x_0)$ 保持定号，对应严格极小值或严格极大值；当 $n$ 为奇数时，$h^n$ 在两侧异号，故 $f(x_0+h)-f(x_0)$ 在左右两侧取相反符号，此时 $x_0$ 不是极值点。
\end{proof}

\section{高阶判别的应用与边界}
在具体计算多项式或光滑函数的极值时，不必真的逐阶求出高阶导数的完整表达式，往往只需写出函数在驻点附近的局部主部，即可直接套用高阶判别法。

\begin{example}
求 $f(x)=x^4(x-1)^3$ 的极值。
\end{example}

\begin{proof}[解]
求导得 $f'(x)=x^3(x-1)^2(7x-4)$，故驻点为 $x=0$、$x=4/7$ 与 $x=1$。分别考察这三点：
\begin{itemize}
\item 在 $x=4/7$ 处，对 $f'$ 求导并代入 $x=4/7$ 时，凡含因子 $7x-4$ 的项均为零，仅对 $7x-4$ 求导所得的项保留，因此 $f''(4/7)=7(4/7)^3(4/7-1)^2>0$，该点为严格极小值点，极小值为 $-(4/7)^4(3/7)^3$。
\item 在 $x=0$ 处，$f(0)=0$，且当 $h\to0$ 时有 $f(h)=h^4(-1+O(h))$。展开式的首个非零项为 $-h^4$，次数为偶数且系数为负，故 $x=0$ 为严格极大值点，极大值为 $0$。
\item 在 $x=1$ 处，$f(1)=0$，且 $f(1+h)=h^3(1+O(h))$。首个非零项为 $h^3$，在 $h$ 的正负两侧异号，故 $x=1$ 不是极值点。
\end{itemize}
\end{proof}

需要明白的是，三个充分条件都不能覆盖所有情形。当函数在驻点处的各阶导数全为零时，高阶判别法会全部失效，此时往往需要回到极值的定义直接比较函数值。例如令 $f(0)=0$，当 $x\ne0$ 时 $f(x)=e^{-1/x^2}$。对任意整数 $m\ge0$，作代换 $t=1/x^2$ 可知，因 $e^t$ 增长快于 $t^{m/2}$，故有 $|x|^{-m}e^{-1/x^2}\to0$（$x\to0$）。又在 $x\ne0$ 处，逐次求导所得的 $f^{(k)}(x)$ 都是 $e^{-1/x^2}$ 与 $1/x$ 的多项式之积，从而在原点处逐阶取差商极限均有 $f^{(k)}(0)=0$。尽管所有阶导数在原点均为零，但对一切 $x\ne0$ 显然有 $f(x)>0=f(0)$，因此原点依然是严格极小值点。

下面的例子进一步表明，甚至可能出现第一、第二充分条件同时失效的情形。

\begin{example}[习题 6.4 第 2 题]
令 $f(0)=0$，当 $x\ne0$ 时 $f(x)=x^4\sin^2(1/x)$。判断原点是否为极值点，并考察第一、第二充分条件能否作出判别。
\end{example}

\begin{proof}[解]
首先直接依据定义判断：对一切实数 $x$ 均有 $f(x)\ge0=f(0)$，故原点是极小值点。又因为对所有正整数 $k$ 都有 $f\bigl(1/(k\pi)\bigr)=0$，在原点的任意小邻域内都存在其他点取到相同的函数值 $0$，所以该极小值不严格。

再检验两个充分条件。由差商 $f(x)/x=x^3\sin^2(1/x)\to0$ 可知 $f'(0)=0$。当 $x\ne0$ 时，
\[
 f'(x)=4x^3\sin^2(1/x)-x^2\sin(2/x).
\]
沿点列 $2/x=\pi/2+2k\pi$ 取值，得 $f'(x)=x^2(2x-1)$，当 $k$ 充分大时为负；沿点列 $2/x=3\pi/2+2k\pi$ 取值，得 $f'(x)=x^2(2x+1)>0$。这两列点均从右侧趋于 $0$，说明 $f'$ 在原点的任意右邻域内都不保持定号，故第一充分条件不能使用。此外，由 $f'(x)/x=4x^2\sin^2(1/x)-x\sin(2/x)\to0$ 可知 $f''(0)=0$，第二充分条件同样失效。在此例中，直接考察函数值的非负性反而最为简洁有效。
\end{proof}

\section{闭区间上的最大值与最小值}
由闭区间上连续函数的最值定理可知，若 $f$ 在 $[a,b]$ 上连续，则它在该区间上必能取得最大值与最小值。若最值在某内点处取得且函数在该点可导，根据费马定理该点必为驻点。因此，求闭区间上连续函数的最值时，无需对每个内点分别判别极值性质，只需将\textbf{区间两端点、内部驻点以及内部不可导点}处的函数值全部列出并加以比较，即可确定最大值与最小值。

\begin{example}
求 $f(x)=\sin x+\cos x$ 在 $[0,2\pi]$ 上的内部极值与区间最值。
\end{example}

\begin{proof}[解]
由 $f'(x)=\cos x-\sin x=0$，在 $(0,2\pi)$ 内解得驻点 $x=\pi/4$ 与 $x=5\pi/4$。计算二阶导数 $f''(x)=-\sin x-\cos x$，得 $f''(\pi/4)=-\sqrt2<0$、$f''(5\pi/4)=\sqrt2>0$，因此 $\pi/4$ 为严格极大值点，$5\pi/4$ 为严格极小值点。再将两驻点与两端点的函数值 $f(0)=f(2\pi)=1$ 一并比较，可知区间最大值为 $f(\pi/4)=\sqrt2$，最小值为 $f(5\pi/4)=-\sqrt2$。
\end{proof}

\begin{example}
求 $f(x)=|2x^3-9x^2+12x|$ 在 $\left[-\frac14,\frac52\right]$ 上的最大值和最小值。
\end{example}

\begin{proof}[解]
将绝对值内的多项式分解为 $x(2x^2-9x+12)$。由于二次因子的判别式为 $81-96<0$ 且首项系数为正，它恒大于零，因此绝对值内的多项式与 $x$ 同号。据此分段求导得
\[
 f'(x)=
 \begin{cases}
 -6(x-1)(x-2), & -\frac14<x<0,\\
 \phantom{-}6(x-1)(x-2), & 0<x<\frac52.
 \end{cases}
\]
在左段 $\left(-\frac14,0\right)$ 内没有驻点，在右段 $\left(0,\frac52\right)$ 内有驻点 $x=1$ 与 $x=2$。而在 $x=0$ 处，左、右差商的极限分别为 $-12$ 与 $12$，故 $x=0$ 是区间内的不可导点。连同两个端点，列出全部五个候选点的函数值：
\[
 \begin{array}{c|ccccc}
 x & -\frac14 & 0 & 1 & 2 & \frac52 \\ \hline
 f(x) & \frac{115}{32} & 0 & 5 & 4 & 5
 \end{array}
\]
比较可知，最小值 $0$ 在不可导点 $x=0$ 处取得，最大值 $5$ 同时在内部驻点 $x=1$ 与右端点 $x=\frac52$ 处取得。
\end{proof}

\section{从变量范围到全局最值}
当定义域不是闭区间时，连续函数未必能取得最值。例如 $f(x)=x$ 在开区间 $(0,1)$ 上的下确界为 $0$、上确界为 $1$，但在区间内部均取不到。不过，若能证明函数在某一点左侧始终递减、右侧始终递增，该点处的函数值便自然成为全局最小值。更一般地，只要连续函数在区间内仅有一个极值点，该极值就必然是全局最值。

\begin{proposition}[习题 6.4 第 5 题]
设 $f$ 在区间 $I$ 上连续，且在 $I$ 内仅有唯一的内点 $x_0$ 为极值点。若 $x_0$ 为极小值点，则 $f(x_0)$ 是 $f$ 在 $I$ 上的全局最小值，且仅在 $x_0$ 处取得。极大值情形同理。
\end{proposition}

\begin{proof}
由于 $x_0$ 是极小值点，在其某个邻域内恒有 $f(x)\ge f(x_0)$。如果该邻域内另有一内点取到与 $f(x_0)$ 相同的值，则该点本身也会成为极小值点，与极值点的唯一性矛盾；因此存在 $x_0$ 的去心邻域，使其中每一点都满足 $f(x)>f(x_0)$。

现假设存在 $y\in I$ 且 $y>x_0$ 使得 $f(y)\le f(x_0)$。由上一步可在 $(x_0,y)$ 内取一点 $z$ 满足 $f(z)>f(x_0)$。因为 $f$ 在闭区间 $[x_0,y]$ 上连续，其最大值必在 $[x_0,y]$ 上取得；而两端点处的函数值均不超过 $f(x_0)<f(z)$，故最大值只能在开区间 $(x_0,y)$ 的某内点处取得。这就产生了异于 $x_0$ 的第二个极值点，矛盾。对 $y<x_0$ 的情形同理可导出矛盾。由此可知，对一切 $y\in I\setminus\{x_0\}$ 均有 $f(y)>f(x_0)$。
\end{proof}

在实际应用问题中，建立目标函数后，还需根据实际含义确定自变量的取值范围，再结合全区间单调性或闭区间连续延拓来确定全局最值。

\begin{example}
一艘轮船每小时的燃料费与船速的立方成正比，当船速为 $10\,\mathrm{km/h}$ 时每小时燃料费为 $6$ 元，其余固定费用每小时为 $96$ 元。船速取多少时，每航行 $1\,\mathrm{km}$ 的总费用最少？
\end{example}

\begin{proof}[解]
设船速为 $x\,\mathrm{km/h}$（$x>0$），则航行 $1\,\mathrm{km}$ 耗时 $1/x$ 小时。由题设知燃料费的比例系数为 $6/10^3=0.006$，因而每千米的总费用及其导数分别为
\[
 C(x)=\frac{0.006x^3+96}{x}=0.006x^2+\frac{96}{x},
 \qquad C'(x)=\frac{0.012(x^3-8000)}{x^2}.
\]
显然，当 $x\in(0,20)$ 时 $C'(x)<0$，当 $x\in(20,+\infty)$ 时 $C'(x)>0$，故 $C(x)$ 在 $(0,20]$ 上严格递减、在 $[20,+\infty)$ 上严格递增，对所有 $x>0$ 恒有 $C(x)\ge C(20)$。因此船速为 $20\,\mathrm{km/h}$ 时每千米总费用最低，最低费用为 $C(20)=7.2$ 元。
\end{proof}

\begin{example}
从边长为 $a>0$ 的正方形纸片四角分别剪去同样大小的小正方形，沿余下的边折成无盖方盒。剪去的小正方形边长取多少时，盒子的容积最大？
\end{example}

\begin{proof}[解]
设剪去的小正方形边长为 $x$，则盒高为 $x$，底面正方形边长为 $a-2x$。由几何可行性可知自变量满足 $0<x<a/2$，容积函数及其导数为
\[
 V(x)=x(a-2x)^2,\qquad V'(x)=(a-2x)(a-6x).
\]
将 $V(x)$ 连续延拓至闭区间 $[0,a/2]$，两端点处均有 $V(0)=V(a/2)=0$，而在开区间 $(0,a/2)$ 内仅有唯一驻点 $x=a/6$，对应容积为 $V(a/6)=2a^3/27>0$。因此 $V(a/6)$ 既是闭区间 $[0,a/2]$ 上的最大值，也是可行范围 $(0,a/2)$ 内所能取得的最大容积。
\end{proof}

\end{document}
